\begin{center}
\hypertarget{native-035-article017}{}\pdfbookmark[2]{17.}{article-017}17.
\end{center}
De indagatione \emph{aequationis centri maximae} pauca adiecisse haud poenitebit. Primo sponte obvium est, differentiam inter anomaliam excentricam et mediam maximum esse pro $E=90^0$, ubi fit $=e$ (in gradibus etc. exprimenda); radius vector in hoc puncto est $=a$, unde $v=90^0+\varphi$, adeoque aequatio centri tota $=\varphi+e$, quae tamen hic non est maximum, quoniam differentia inter $v$ et $E$ adhuc ultra $\varphi$ crescere potest. \emph{Haecce} differentia fit maximum pro $\mathrm{d}(v-E)=0$ sive pro $\mathrm{d}v=\mathrm{d}E$, ubi excentricitas manifesto ut constans spectanda est. Qua suppositione quum generaliter fiat $\frac{\mathrm{d}v}{\sin v}=\frac{\mathrm{d}E}{\sin E}$, patet, in eo puncto, ubi differentia inter $v$ et $E$ maximum est, esse debere $\sin v=\sin E$; unde erit, per aequatt. VIII, III,
\[
r=a\cos\varphi
\]
\[
e\cos E=1-\cos\varphi
\]
sive
\[
\cos E=+\operatorname{tang}\frac{1}{2}\varphi.
\]
Perinde invenitur $\cos v=-\operatorname{tang}\frac{1}{2}\varphi$, quapropter erit\footnote{Ad ea maxima, quae inter aphelium et perihelium iacent, non opus est respicere, quum manifesto ab iis, quae inter perihelium et aphelium sita sunt, in signis tantum differant.}
\[
v=90^0+\operatorname{arc.}\sin\operatorname{tang}\frac{1}{2}\varphi
\]
\[
E=90^0-\operatorname{arc.}\sin\operatorname{tang}\frac{1}{2}\varphi;
\]
hinc porro $\sin E=\BareRoot{(1-\operatorname{tang}\frac{1}{2}\varphi^2)}=\frac{\BareRoot{\cos\varphi}}{\cos\frac{1}{2}\varphi}$, ita ut aequatio centri tota in hoc puncto fiat
\[
=2\operatorname{arc.}\sin\operatorname{tang}\frac{1}{2}\varphi+2\sin\frac{1}{2}\varphi\cdot\BareRoot{\cos\varphi}
\]
parte secunda in gradibus etc. expressa.

In eo denique puncto, ubi tota aequatio centri ipsa maximum est, fieri debet $\mathrm{d}v=\mathrm{d}M$, adeoque, secundum art. 15, $r=a\BareRoot{\cos\varphi}$; hinc fit
\SourceLeaf{036}
\[
\cos v=-\frac{1-\cos\varphi^{\frac{3}{2}}}{e}
\]
\[
\cos E=\frac{1-\BareRoot{\cos\varphi}}{e}=\frac{1-\cos\varphi}{e(1+\BareRoot{\cos\varphi})}=\frac{\operatorname{tang}\frac{1}{2}\varphi}{1+\BareRoot{\cos\varphi}},
\]
per quam formulam $E$ ultima praecisione determinare licet. Inventa $E$, erit per aequ. X, XII aequatio centri $=2\operatorname{arc.}\sin\frac{\sin\frac{1}{2}\varphi\sin E}{\BareRoot[4]{\cos\varphi}}+e\sin E$\HistoricalNoteEight{Handschriftliche Bemerkung:] $\sin E=\BareRoot{\frac{\cos\varphi+2\BareRoot{\cos\varphi}+\cos\varphi^{\frac{3}{2}}}{(1+\BareRoot{\cos\varphi})(1+\cos\varphi)}},$ Aequatio centri $=2\operatorname{arc.}\sin\BareRoot{\left(\frac{1}{1+\cos\varphi}-\frac{\BareRoot[4]{(1-e e)}}{2}\right)}+e\sin E=\operatorname{arc.}\cos\frac{-1+\BareRoot{\cos\varphi}+\cos\varphi+\cos\varphi^{\frac{3}{2}}}{1+\cos\varphi}+e\sin E=\operatorname{arc.}\sin(\operatorname{tg}\frac{1}{2}\varphi^2\operatorname{tg}E)+e\sin E.$}. Expressioni aequationis centri maximae per seriem secundum potestates excentricitatis progredientem, quam plures auctores tradiderunt, hic non immoramur. Ut exemplum habeatur, conspectum trium maximorum, quae hic contemplati sumus, pro Iunone adiungimus, ubi excentricitas secundum elementa novissima $=0{,}255\,4996$ supposita est.
\begin{center}
\footnotesize
\setlength{\tabcolsep}{3pt}
\renewcommand{\arraystretch}{1.3}
\begin{tabular}{|c|c|c|c|c|}
\hline
Maximum & $E\quad\text{·}$ & $E-M$ & $v-E$ & $v-M$\\
\hline
$E-M$ & $90^0\ 0^{\prime}\ 0^{\prime\prime}$ & $14^0\,38^{\prime}\,20^{\prime\prime}{,}57$ & $14^0\,48^{\prime}\,11^{\prime\prime}{,}48$ & $29^0\,26^{\prime}\,32^{\prime\prime}{,}05$\\
$v-E$ & $82\ \ 32\ \ 9$ & $14\ \ 30\ \ 54{,}01$ & $14\ \ 55\ \ 41{,}79$ & $29\ \ 26\ \ 35{,}80$\\
$v-M$ & $86\ \ 14\ \ 40$ & $14\ \ 36\ \ 27{,}39$ & $14\ \ 53\ \ 49{,}57$ & $29\ \ 30\ \ 16{,}96$\\
\hline
\end{tabular}
\end{center}
